\chapter{Conclusions} \label{chapter:6}
This project investigated the feasibility and performance of LEO satellite constellations as a broadband connectivity solution for rural Ireland. Using orbital propagation and coverage modelling, the performance of various constellations were evaluated in terms of coverage availability and downlink capacity. These results were then compared to terrestrial 5G networks. 

\section{Key Findings}
The results demonstrate that constellation design has a major impact on coverage performance over Ireland. Starlink cosnsitently achieved the highest covereage availability, with almost continuous coverage in parts of the country. This performance can largely be attributed to Starlink's high satellite density, low orbital altitdues and otbital shells which reguarly pass over parts of Ireland.

OneWeb showed signifcantly lower coverage availability, with satellites passing infrequently and with short coverage windows. Coverage was slightly superior in Northern regions but still only reached a maximum coverage availability of $11.67\%$. However, when coverage was available a strong data rate was calculated with a peak of 240.89 Mbps, showing that it is capable of delivering 5G level broadband. 

Kuiper also lacked in terms of coverage availability, with only a max coverage of $7.88\%$ in the southern parts of the country, however, of the three constellations modelled, it had the weakest availability and the weakest coverage speeds, at only a maximum data rate of 207.71 Mbps. 

While all three constellations are capable of providing 5G level coverage, the poor availability in Ireland from Kuiper and OneWeb (and in some parts of the country, Starlink), make them an unviable option for rural connectivity as they stand. 

\section{Comparison with Terrestrial 5G}
The analysis of terrestrial 5G coverage shows a dependence on population density in terms of infrastructure investment. High levels of coverage were more likely to be seen in urban regions, while significant gaps could be seen in rural areas. Where 5G exists, it is often capable of a stronger connection than LEO satellites, however, coverage is not uniform nationwide. 

LEO constellations however, provide broad coverage independent of local poplation density and infrastructure. This supports the conlcusion that LEO satellite broadband should be used as a complementary technology rather than a replacement for terrestrial netowrks. In urban areas, terreistral 5G remains optimal, while LEO systems offer a viable alternative for rural areas.  

\section{Limitations of Study}
This study employs a simplified model focused on coverage availability and downlink capacity. Layered network performance, uplink performance, latency variability, and effects of various levels of trafic are not modelled. On top of this, constellations paramaters were based on publically available informatin rather than operation performance data. These simplifications limit the accuracyt of the performance estimates, but they are appropiate for a comparative, system level comparison of coverage over a region. 



\section{Project Reflection}
This section reflects on the execution of the project, evaluating the planning decisions, technical approach, challenges encountered, and lessons learned throughout the development process. 

\subsection{Technical Approach and Design Choices}

\subsection{Challenges Encountered}

\subsection{Evaluation of outcomes}

\subsection{Lessons Learned and Future Improvements}